% Generated draft: requires mathematical and source review.
\documentclass[11pt]{article}
\usepackage[margin=25mm]{geometry}
\usepackage{fontspec}
\setmainfont{Latin Modern Roman}
\usepackage{amsmath,amssymb,mathtools,unicode-math}
\setmathfont{Latin Modern Math}
\usepackage{enumitem,needspace,xurl,hyperref}
\setlength{\parindent}{0pt}
\setlength{\parskip}{5pt}
\setlength{\emergencystretch}{3em}
\newcommand{\N}{\mathbb{N}}
\newcommand{\Z}{\mathbb{Z}}
\newcommand{\Q}{\mathbb{Q}}
\newcommand{\R}{\mathbb{R}}
\newcommand{\C}{\mathbb{C}}
\newcommand{\F}{\mathbb{F}}
\newcommand{\A}{\mathbb{A}}
\newcommand{\PP}{\mathbb{P}}
\newcommand{\E}{\mathbb{E}}
\newcommand{\eps}{\varepsilon}
\newcommand{\dd}{\,\mathrm{d}}
\begin{document}
\section*{UnsolvedMath ID: 10400132}
\textbf{Source problem code:} catalogue entry
\subsection{Definitions and mathematical statement}
T. Ohtsuki's Problem 7.17 seeks a ``complex structure'' on the set of \(3\)-manifolds, in the following sense. Let \(\mathfrak M_3\) be the source's set of manifold equivalence classes, with equivalence and manifold-category conventions still unspecified. Seek a single complex variety \(X\) and an embedding, or alternatively an immersion,
\[
 \iota:\mathfrak M_3\longrightarrow X,
\]
whose values on sufficiently long Dehn fillings of every relevant knot complement extend holomorphically in the complex parameter referred to in the preceding source text.

In the usual surgery interpretation of the notation, if \(K\subset N\) is a knot and a peripheral basis \((\mu,\lambda)\) has been chosen on the boundary of \(N\setminus\operatorname{int}\nu K\), then \(N_{K;(p,q)}\) is its filling along slope \(p\mu+q\lambda\). Primitive pairs and the identification of opposite pairs must be specified in that interpretation; they are not stated in the extract. The required restriction is to
\[
 \{N_{K;(p,q)}:p^2+q^2\gg0\}.
\]
For any \((\text{hyperbolic})\) knot \(K\) in any \(3\)-manifold \(N\), the map assigning \(\iota(N_{K;(p,q)})\) must extend to a holomorphic map in the source's ``above mentioned complex parameter.'' The parameter, its domain, the corresponding parameter values of the fillings, and the meaning and dependence of the sufficiently-large threshold are all missing. Hence no specific interpolation formula or complex coordinate is imposed here.

An embedding or immersion of a bare set has no differential or complex-geometric meaning until a source structure and the intended definition are supplied; set-theoretic injectivity alone does not express the requested analytic compatibility. The parenthetical ``hyperbolic'' also leaves the knot scope to source review. The excerpt presents this as a hoped-for framework for studying manifolds and invariants and cross-refers to Problems 7.17 and 10.16, without a further mathematical criterion.
\subsection{Short English statement}
Place the set of three-manifolds in a common complex variety so that sufficiently long Dehn-filling families interpolate holomorphically in the intended complex surgery parameter. The parameter and the meaning of embedding or immersion for the source set are absent, so the proposed framework is not yet a precise existence theorem.
\subsection{Sources}
\begin{itemize}
\item Local UnsolvedMath registry, permanent ID \(10400132\), exact source code \texttt{catalogue entry}, Problem 7.17, attributed to T. Ohtsuki. The title is truncated after ``(T.'' The supplied surrounding comment cross-refers to Problems 7.17 and 10.16, including the apparent self-reference.
\item Ohtsuki, \emph{Problems on invariants of knots and 3-manifolds} (2002), Problem 7.17, PDF page 112: \url{https://msp.org/gtm/2002/04/gtm-2002-04-024s.pdf}. Bold printed headings distinguish 231 numbered items from surrounding references.
\item Extraction pdf-font-structure, accessed 2026-07-29. Source evidence NEEDS\_REVIEW presents an open problem and requires a current primary-literature audit. Supplied classification OPEN-TRIAGE calls this complex-structure program not obviously resolved; no additional primary reference is supplied.
\item Dataset attribution: Ulam AI / UnsolvedMath Contributors, CC BY 4.0. Source rights NEEDS\_REVIEW: publicly accessible author/publisher copy, redistribution terms requiring review.
\end{itemize}
\end{document}
